\documentclass[11pt]{article}

\usepackage[margin=1in]{geometry}
\usepackage{amsmath, amssymb, amsthm}
\usepackage{mathtools}
\usepackage{algorithm}
\usepackage{algpseudocode}
\usepackage[colorlinks=true, linkcolor=black, citecolor=black, urlcolor=blue]{hyperref}

%--------------------------------------------------------------
% Theorem environments
%--------------------------------------------------------------
\theoremstyle{plain}
\newtheorem{theorem}{Theorem}
\newtheorem{lemma}{Lemma}
\newtheorem{proposition}{Proposition}
\newtheorem{corollary}{Corollary}

\theoremstyle{definition}
\newtheorem{definition}{Definition}
\newtheorem{construction}{Construction}
\newtheorem{remark}{Remark}

%--------------------------------------------------------------
% Notation macros
%--------------------------------------------------------------
\newcommand{\bit}{\{0,1\}}
\newcommand{\secp}{n}                       % security parameter
\newcommand{\adv}{\mathcal{A}}
\newcommand{\dist}{\mathcal{D}}
\newcommand{\Enc}{\mathsf{Enc}}
\newcommand{\Dec}{\mathsf{Dec}}
\newcommand{\Gen}{\mathsf{Gen}}
\newcommand{\Real}{\mathsf{real}}
\newcommand{\negl}{\mathsf{negl}}
\newcommand{\PPT}{\mathrm{PPT}}
\newcommand{\uecho}{\ensuremath{\leftarrow}}
\DeclarePairedDelimiter{\abs}{\lvert}{\rvert}
\newcommand{\getsr}{\overset{\$}{\leftarrow}}   % uniform sampling
\newcommand{\Adv}{\mathbf{Adv}}
\newcommand{\PrivK}{\mathsf{PrivK}}
\newcommand{\eav}{\mathrm{eav}}
\newcommand{\prg}{\mathrm{prg}}
\newcommand{\Pron}{\Pi}
\newcommand{\Ptilde}{\widetilde{\Pi}}

\allowdisplaybreaks

\title{\bf Security of the Stream Cipher \\ from Pseudorandomness of its Generator}
\author{}
\date{\today}

\begin{document}
\maketitle

\begin{abstract}
We give a self-contained, reductionist proof that the canonical stream-cipher
encryption scheme---encryption by XOR with the output of a pseudorandom
generator (PRG)---achieves indistinguishability of encryptions in the presence
of an eavesdropper, provided the underlying generator is a secure PRG. The
argument is presented in prose, with the single reduction algorithm displayed as
pseudocode. The reduction is advantage-preserving: the eavesdropping advantage
against the cipher is bounded above by the distinguishing advantage against the
generator, with no loss. We isolate the one-time-pad as the information-theoretic
anchor of the hybrid and remark on tightness, on the necessity of the
assumption, and on the failure of the single-message guarantee to extend to the
multi-message setting for the stateless scheme.
\end{abstract}

%==============================================================
\section{Preliminaries and notation}
%==============================================================

We work in the asymptotic (uniform-complexity) model of computation. Throughout,
$\secp \in \mathbb{N}$ denotes the security parameter, and all algorithms receive
$1^\secp$ as (implicit or explicit) input. We write $x \getsr S$ for the
assignment to $x$ of an element drawn uniformly at random from a finite set $S$,
and $y \leftarrow \adv(x)$ for the assignment to $y$ of the (possibly
randomized) output of an algorithm $\adv$ on input $x$. By $\PPT$ we mean
probabilistic polynomial-time (in $\secp$).

\begin{definition}[Negligible function]
A function $\negl \colon \mathbb{N} \to \mathbb{R}_{\geq 0}$ is \emph{negligible}
if for every positive polynomial $p$ there exists $N \in \mathbb{N}$ such that
$\negl(\secp) < 1/p(\secp)$ for all $\secp > N$. Equivalently,
$\negl(\secp) = \secp^{-\omega(1)}$.
\end{definition}

The class of negligible functions is closed under addition and under
multiplication by any fixed polynomial; we use both facts without comment.

For a bit string $x$ we write $\abs{x}$ for its length and $\oplus$ for bitwise
exclusive-or; for equal-length strings $x,y$ the map $z \mapsto z \oplus x$ is a
bijection on $\bit^{\abs{x}}$, a fact we invoke repeatedly.

%==============================================================
\section{The two definitions in play}
%==============================================================

\subsection{Pseudorandom generators}

Let $\ell \colon \mathbb{N} \to \mathbb{N}$ be a polynomial with
$\ell(\secp) > \secp$ for all $\secp$; $\ell$ is the \emph{expansion factor}.

\begin{definition}[Pseudorandom generator]\label{def:prg}
A deterministic polynomial-time computable function
$G \colon \bit^{*} \to \bit^{*}$ with $\abs{G(s)} = \ell(\abs{s})$ for all $s$ is
a \emph{pseudorandom generator} with expansion $\ell$ if for every $\PPT$
distinguisher $\dist$ the \emph{distinguishing advantage}
\[
  \Adv^{\prg}_{G,\dist}(\secp)
  \;\coloneqq\;
  \Bigl\lvert
     \Pr_{s \getsr \bit^{\secp}}\!\bigl[\dist(G(s)) = 1\bigr]
     \;-\;
     \Pr_{r \getsr \bit^{\ell(\secp)}}\!\bigl[\dist(r) = 1\bigr]
  \Bigr\rvert
\]
is negligible in $\secp$. Both probabilities are additionally over the internal
randomness of $\dist$.
\end{definition}

The seed length $\secp$ is exactly the security parameter, so a uniform seed is
an $\secp$-bit string and its image is an $\ell(\secp)$-bit string; the gap
$\ell(\secp) - \secp$ is the number of bits by which $G$ ``stretches'' its input.
It is precisely this stretch that makes the definition nontrivial: were
$\ell(\secp) = \secp$, the identity function would satisfy
Definition~\ref{def:prg}.

\subsection{Private-key encryption and eavesdropping security}

A \emph{private-key encryption scheme} is a triple
$\Pron = (\Gen, \Enc, \Dec)$ of $\PPT$ algorithms with an associated message
space $\{\mathcal{M}_\secp\}_\secp$, satisfying the usual correctness condition:
for every $\secp$, every key $k$ in the support of $\Gen(1^\secp)$, and every
$m \in \mathcal{M}_\secp$, we have $\Dec_k(\Enc_k(m)) = m$ with probability $1$.

We adopt the standard indistinguishability formulation of eavesdropping
security. For a scheme $\Pron$, an adversary $\adv$, and a value $\secp$, define
the experiment $\PrivK^{\eav}_{\adv,\Pron}(\secp)$ as follows.

\begin{enumerate}
  \item $\adv(1^\secp)$ outputs a pair of messages
        $m_0, m_1 \in \mathcal{M}_\secp$ with $\abs{m_0} = \abs{m_1}$.
  \item A key is generated, $k \leftarrow \Gen(1^\secp)$, a challenge bit is
        chosen, $b \getsr \bit$, and the \emph{challenge ciphertext}
        $c \leftarrow \Enc_k(m_b)$ is computed and handed to $\adv$.
  \item $\adv$ outputs a bit $b'$. The experiment \emph{evaluates to $1$}
        (``$\adv$ succeeds'') if $b' = b$, and to $0$ otherwise.
\end{enumerate}

\begin{definition}[Indistinguishable encryptions under eavesdropping]
\label{def:eav}
$\Pron$ has \emph{indistinguishable encryptions in the presence of an
eavesdropper}, or is \emph{$\eav$-secure}, if for every $\PPT$ adversary $\adv$
the \emph{eavesdropping advantage}
\[
  \Adv^{\eav}_{\Pron,\adv}(\secp)
  \;\coloneqq\;
  \Pr\!\bigl[\PrivK^{\eav}_{\adv,\Pron}(\secp) = 1\bigr] - \tfrac{1}{2}
\]
is negligible in $\secp$.
\end{definition}

The subtracted $1/2$ is the success probability of the trivial adversary that
ignores $c$ and guesses; $\eav$-security asks that no efficient strategy beat
this baseline by a non-negligible margin. We record for later use the
information-theoretic gold standard against which the computational notion is
calibrated.

\begin{definition}[Perfect secrecy]\label{def:perfect}
$\Pron$ is \emph{perfectly secret} if for every $\secp$, every adversary $\adv$
whatsoever (of arbitrary computational power), and the experiment above,
$\Pr[\PrivK^{\eav}_{\adv,\Pron}(\secp)=1] = \tfrac12$ exactly.
\end{definition}

The relevant instance is Vernam's one-time pad, whose perfect secrecy we invoke
as a lemma.

\begin{lemma}[One-time pad]\label{lem:otp}
Fix $\secp$ and let $L \coloneqq \ell(\secp)$. Let $\Ptilde$ denote the scheme
with message space $\bit^{L}$ in which $\Gen$ outputs a uniform pad
$p \getsr \bit^{L}$ and $\Enc_p(m) = p \oplus m$. Then $\Ptilde$ is perfectly
secret; in particular
$\Pr[\PrivK^{\eav}_{\adv,\Ptilde}(\secp) = 1] = \tfrac{1}{2}$ for every $\adv$.
\end{lemma}

\begin{proof}
Condition on any fixed messages $m_0, m_1 \in \bit^{L}$ output by $\adv$ and on
any fixed challenge bit $b$. Because $p$ is uniform on $\bit^{L}$ and independent
of $b$, and $x \mapsto x \oplus m_b$ is a bijection of $\bit^{L}$, the challenge
$c = p \oplus m_b$ is uniform on $\bit^{L}$; crucially, its distribution does not
depend on $b$. Hence the view of $\adv$ is statistically independent of $b$, and
$\adv$'s output $b'$ is independent of $b$. Averaging over the uniform choice of
$b \in \bit$ gives $\Pr[b' = b] = \tfrac12$.
\end{proof}

%==============================================================
\section{The construction and the theorem}
%==============================================================

\begin{construction}[Fixed-length stream cipher]\label{con:cipher}
Let $G$ be a deterministic polynomial-time function with expansion $\ell$. Define
$\Pron = (\Gen, \Enc, \Dec)$ over the message space
$\mathcal{M}_\secp = \bit^{\ell(\secp)}$ by
\[
  \Gen(1^\secp)\colon\ k \getsr \bit^{\secp};
  \qquad
  \Enc_k(m) = G(k) \oplus m;
  \qquad
  \Dec_k(c) = G(k) \oplus c .
\]
\end{construction}

Correctness is immediate:
$\Dec_k(\Enc_k(m)) = G(k) \oplus (G(k) \oplus m) = m$. The scheme is the standard
(synchronous, fixed-output-length) stream cipher: the generator expands a short
key into a long \emph{keystream} $G(k)$, and encryption is XOR of the plaintext
with the keystream, exactly as in the one-time pad but with a pseudorandom pad in
place of a truly random one. The key is $\secp$ bits while the message is
$\ell(\secp) > \secp$ bits, so---unlike the one-time pad---the scheme encrypts
messages longer than its key.

\begin{theorem}\label{thm:main}
If $G$ is a pseudorandom generator with expansion $\ell$, then the scheme
$\Pron$ of Construction~\ref{con:cipher} has indistinguishable encryptions in the
presence of an eavesdropper (Definition~\ref{def:eav}). Quantitatively, for every
$\PPT$ eavesdropper $\adv$ there is a $\PPT$ distinguisher $\dist$, running in
time $t_\adv + O(\ell(\secp))$ where $t_\adv$ is the running time of $\adv$, such
that
\begin{equation}\label{eq:tight}
  \Adv^{\eav}_{\Pron,\adv}(\secp)
  \;\leq\;
  \Adv^{\prg}_{G,\dist}(\secp)
  \qquad\text{for all } \secp .
\end{equation}
\end{theorem}

The asymptotic statement follows from the concrete
bound~\eqref{eq:tight}: the right-hand side is negligible because $G$ is a PRG,
hence so is the left-hand side. We therefore prove the quantitative form.

%==============================================================
\section{Proof of Theorem~\ref{thm:main}}
%==============================================================

\paragraph{Strategy.} We exhibit a single distinguisher $\dist$ that turns any
eavesdropper $\adv$ into an attack on $G$. The distinguisher receives a
challenge string $w \in \bit^{\ell(\secp)}$, promised to be either $G(s)$ for a
uniform seed $s$ or a uniform string $r$, and it must decide which. It does so by
\emph{using $w$ as the keystream} in a simulated run of the eavesdropping
experiment against $\adv$, and betting that $w$ is pseudorandom exactly when
$\adv$ wins the simulated experiment. The two horns of the PRG promise then
correspond to the two schemes $\Pron$ (real keystream) and $\Ptilde$ (random
pad, i.e.\ the one-time pad of Lemma~\ref{lem:otp}), and the reduction reads off
the theorem by comparing $\adv$'s success in each.

\begin{algorithm}[t]
\caption{Distinguisher $\dist$ against $G$, with black-box access to eavesdropper $\adv$}
\label{alg:dist}
\begin{algorithmic}[1]
\Require challenge string $w \in \bit^{\ell(\secp)}$ (equal to $G(s)$ or to uniform $r$)
\State $(m_0, m_1) \leftarrow \adv(1^\secp)$
       \Comment{obtain the challenge messages, $\abs{m_0}=\abs{m_1}=\ell(\secp)$}
\State $b \getsr \bit$
       \Comment{internal challenge bit, chosen by $\dist$}
\State $c \gets w \oplus m_b$
       \Comment{use the challenge string \emph{as the keystream}}
\State $b' \leftarrow \adv(c)$
       \Comment{deliver the simulated challenge ciphertext}
\If{$b' = b$}
    \State \Return $1$ \Comment{guess ``$w$ is pseudorandom''}
\Else
    \State \Return $0$ \Comment{guess ``$w$ is uniform''}
\EndIf
\end{algorithmic}
\end{algorithm}

The distinguisher is displayed as Algorithm~\ref{alg:dist}. It is a $\PPT$
machine: it runs $\adv$ once, performs one XOR of two $\ell(\secp)$-bit strings,
and one bit comparison, so its running time is $t_\adv + O(\ell(\secp))$ as
claimed. We now analyze $\dist$ under each of the two input distributions of
Definition~\ref{def:prg}.

\paragraph{Case 1: $w = G(s)$ for $s \getsr \bit^{\secp}$.} We claim the view
$\dist$ presents to $\adv$ is \emph{identically distributed} to $\adv$'s view in
the real experiment $\PrivK^{\eav}_{\adv,\Pron}(\secp)$. Indeed, in that
experiment the key is $k \getsr \bit^{\secp}$ and the challenge ciphertext is
$\Enc_k(m_b) = G(k) \oplus m_b$; in the simulation the ``key'' is the uniform
seed $s$ underlying $w = G(s)$, and the challenge ciphertext is
$w \oplus m_b = G(s) \oplus m_b$. Since $s$ and $k$ are both uniform on
$\bit^{\secp}$ and $\dist$ chooses $b$ exactly as the experiment does, the joint
distribution of $(m_0, m_1, b, c)$---and hence of $\adv$'s output $b'$---is the
same in both. By construction $\dist$ outputs $1$ iff $b' = b$, i.e.\ iff the
simulated experiment evaluates to $1$. Therefore
\begin{equation}\label{eq:case1}
  \Pr_{s \getsr \bit^{\secp}}\!\bigl[\dist(G(s)) = 1\bigr]
  \;=\;
  \Pr\!\bigl[\PrivK^{\eav}_{\adv,\Pron}(\secp) = 1\bigr] .
\end{equation}

\paragraph{Case 2: $w = r$ for $r \getsr \bit^{\ell(\secp)}$.} Now the challenge
ciphertext handed to $\adv$ is $c = r \oplus m_b$ with $r$ uniform on
$\bit^{\ell(\secp)}$ and independent of $b$. This is exactly the challenge
produced in the eavesdropping experiment against the one-time pad $\Ptilde$ of
Lemma~\ref{lem:otp} (whose message space is $\bit^{\ell(\secp)}$ and whose
``key'' is a uniform pad). Since $\dist$ again outputs $1$ iff $\adv$ succeeds,
\[
  \Pr_{r \getsr \bit^{\ell(\secp)}}\!\bigl[\dist(r) = 1\bigr]
  \;=\;
  \Pr\!\bigl[\PrivK^{\eav}_{\adv,\Ptilde}(\secp) = 1\bigr]
  \;=\; \tfrac{1}{2},
\]
the last equality by Lemma~\ref{lem:otp}. Note this holds for the very same $\adv$
and the very same reduction; the perfect secrecy of $\Ptilde$ pins the second
probability to $1/2$ irrespective of $\adv$'s strategy.

\paragraph{Combining.} Subtracting the two cases and using that the left-hand
sides differ by at most the distinguishing advantage,
\begin{align*}
  \Adv^{\eav}_{\Pron,\adv}(\secp)
  &= \Pr\!\bigl[\PrivK^{\eav}_{\adv,\Pron}(\secp) = 1\bigr] - \tfrac12 \\
  &= \Pr_{s}\!\bigl[\dist(G(s)) = 1\bigr]
     - \Pr_{r}\!\bigl[\dist(r) = 1\bigr]
     && \text{(by \eqref{eq:case1} and Case 2)} \\
  &\leq \Bigl\lvert
        \Pr_{s}\!\bigl[\dist(G(s)) = 1\bigr]
        - \Pr_{r}\!\bigl[\dist(r) = 1\bigr]
        \Bigr\rvert
     \;=\; \Adv^{\prg}_{G,\dist}(\secp).
\end{align*}
This is exactly~\eqref{eq:tight}. Since $G$ is a PRG and $\dist$ is $\PPT$, the
right-hand side is negligible; being nonnegative and bounded above by a
negligible function, $\Adv^{\eav}_{\Pron,\adv}(\secp)$ is negligible. As $\adv$
was an arbitrary $\PPT$ eavesdropper, $\Pron$ is $\eav$-secure. \qed

%==============================================================
\section{Remarks}
%==============================================================

\begin{remark}[The hybrid, made explicit]
The proof is a two-point hybrid argument in disguise. Define
$H_0$ to be $\adv$'s success probability against the real scheme $\Pron$ and
$H_1$ its success probability against the idealized scheme $\Ptilde$ that XORs
with a truly random pad. Lemma~\ref{lem:otp} evaluates the ideal endpoint
exactly, $H_1 = \tfrac12$, using no computational assumption; the single
distinguisher of Algorithm~\ref{alg:dist} bounds the gap $\abs{H_0 - H_1}$ by the
PRG advantage. All computational content is thus localized to the one step where
$G(k)$ is swapped for a uniform string, which is the only place the assumption is
used and the only source of the (negligible) security loss.
\end{remark}

\begin{remark}[Tightness]
The reduction is advantage-preserving and rewinding-free: the bound
\eqref{eq:tight} carries a multiplicative constant of exactly $1$ and additive
loss $0$, and $\dist$ invokes $\adv$ a single time with essentially no overhead
beyond one $\ell(\secp)$-bit XOR. Consequently the cipher inherits the concrete
security of $G$ verbatim: an adversary spending work $t$ and achieving
eavesdropping advantage $\varepsilon$ yields a PRG distinguisher of work
$\approx t$ and advantage $\geq \varepsilon$. No security is lost in passing from
the primitive to the scheme.
\end{remark}

\begin{remark}[On necessity of the assumption]
Sufficiency of pseudorandomness is what the theorem asserts; a converse also
holds for this construction, so the assumption is not an artifact of the proof.
If $G$ fails Definition~\ref{def:prg}, witnessed by a $\PPT$ $\dist$ with
non-negligible $\Adv^{\prg}_{G,\dist}$, then $\Pron$ fails
Definition~\ref{def:eav}: the eavesdropper submits $m_0 = 0^{\ell(\secp)}$ and a
uniformly random $m_1$, and on receiving $c$ runs $\dist$ on $c \oplus m_0 = c$,
guessing $b'=0$ when $\dist$ answers ``pseudorandom.'' When $b=0$ the string fed
to $\dist$ is $G(k)$; when $b=1$ it is $G(k)\oplus m_1$, uniform and independent
of $G(k)$ because $m_1$ is a fresh uniform pad, so it is distributed as a truly
random string. The eavesdropper's advantage is thus $\tfrac12\Adv^{\prg}_{G,\dist}$,
non-negligible. Hence, for this scheme, $\eav$-security and pseudorandomness of
$G$ are equivalent.
\end{remark}

\begin{remark}[Scope: single message only]
Definition~\ref{def:eav} models a \emph{single} ciphertext. The theorem does
\emph{not} certify the stateless scheme of Construction~\ref{con:cipher} under
multiple encryptions with a fixed key: two messages encrypted under the same
$k$ produce $c_1 \oplus c_2 = m_1 \oplus m_2$, leaking the plaintext XOR and
breaking any reasonable multi-message notion---precisely the two-time-pad
failure. Security across many messages is recovered either by making the
generator stateful (advancing the keystream so that disjoint keystream segments
encrypt successive messages, as in a synchronized stream cipher) or by moving to
a randomized/nonce-based construction from a pseudorandom \emph{function}. Each
requires its own definition (e.g.\ multiple-encryption or CPA indistinguishability)
and its own reduction; neither is subsumed by Theorem~\ref{thm:main}.
\end{remark}

\begin{remark}[Uniform vs.\ non-uniform adversaries]
The argument is agnostic to the adversary model. Read with $\PPT$ meaning
uniform probabilistic polynomial time, it establishes security against uniform
eavesdroppers from security of $G$ against uniform distinguishers; read with
$\dist$ and $\adv$ as polynomial-size circuit families and ``negligible'' quantified
accordingly, the identical reduction establishes the non-uniform statement, since
$\dist$ merely wraps $\adv$ with a fixed amount of extra computation and no
additional nonuniform advice.
\end{remark}

\end{document}
